\documentclass[10pt,a4paper]{amsart}
\usepackage[margin=19mm]{geometry}
\usepackage{amsmath,amssymb,mathtools}
\usepackage[scr=rsfs,cal=euler]{mathalfa}
\usepackage{xfrac}
\usepackage[unicode,hidelinks,bookmarks=false]{hyperref}
\newcommand{\ie}{i.e.\ }
\newcommand{\cf}{cf.\ }
\newcommand{\resp}{respectively\ }
\newcommand{\Subsection}{\subsection}
\providecommand{\nobreakdash}{\nobreak\hbox{-}}
\setlength{\emergencystretch}{2em}
\multlinegap0pt
\usepackage[all]{xy}
\usepackage{bm}
\usepackage[shortlabels]{enumitem}
\usepackage{amssymb}
\usepackage{mathtools}
\usepackage{dsfont}
\usepackage{bbm}
\newtheorem{lemma}{Lemma}
\newcommand{\A}{\operatorname{A}}
\newcommand{\F}{\operatorname{F}}
\newcommand{\FP}{\mathbb{Z} \myast \F_k}
\newcommand{\SAut}{\operatorname{SAut}}
\newcommand{\la}{\langle}
\newcommand{\myast}{\hspace{-1pt} \ast \hspace{-1pt}}
\newcommand{\ra}{\rangle}
\title{sampling elements of a finite group: efficiency of the product replacement algorithm with an accumulator.}
\author{Michał Marcinkowski, Piotr Mizerka}
\date{Source: arXiv:2512.18302v1 (CC BY 4.0)}
\begin{document}
\maketitle

\noindent\emph{Source and licence.} sampling elements of a finite group: efficiency of the product replacement algorithm with an accumulator.. Version arXiv:2512.18302v1 (CC BY 4.0), distributed by arXiv under the Creative Commons Attribution 4.0 International licence, http://creativecommons.org/licenses/by/4.0/, as recorded in the arXiv licence metadata for this exact version and shown on the abstract page https://arxiv.org/abs/2512.18302v1. Full licence notice: \texttt{licenses/arxiv-2512.18302-CC-BY-4.0.txt}.\par\smallskip

\noindent\textit{Editorial scope.} Counted unit: the article's lemma lemma (source0.tex lines 612--614). The article's complete original proof of it is delivered with it (source0.tex lines 616--645, one \texttt{proof} environment of 30 lines). No mathematics is written, repaired or re-derived: the statement and the proof are the article's own lines, in the article's own notation, and no retained line is substituted or re-flowed.  No further source-local context is needed: every cross-reference of every retained line resolves inside the two delivered spans.  Nothing was omitted from any retained span, so the package carries no omission record.  No retained line contains a citation, so the wrapper carries no bibliography and no bibliography entry.  No retained line contains a figure environment, an image-inclusion command or drawing code, so no image is delivered and the package declares no asset dependency.  Editorial formatting only, in addition: an AMS \texttt{amsart} preamble that restates the article's own declarations and macros used by the retained lines, namely the environments \texttt{lemma} on separate counters, together with the standard packages those lines need.  The retained environments print in this document as Lemma 1 in order of appearance, whereas the article prints the same items with the numbering of its own full document.  The complete native source of the article as distributed in the arXiv e-print for this exact version is bundled unchanged as \texttt{sources/191375/source0.tex}.
\begin{lemma}
The group $\A_k$ is isomorphic to $\F_k^2 \rtimes \SAut(\F_k)$, where the action of $\SAut(\F_k)$ on $\F_k^2~=~\F_k \times \F_k$ is the direct product of the canonical action. 
\end{lemma}

\begin{proof}
Let $\F_k < \FP$, where $\F_k = \la x_1,\ldots,x_k \ra$.
Note that $N$-generators generate $\SAut(\F_k)$, thus $\SAut(\F_k) < \A_k$.
Moreover, elements of $\A_k$ fix $\F_k$, hence by restricting to this subgroup, we have the retraction:
$$
p \colon \A_k \twoheadrightarrow \SAut(\F_k).
$$

Thus $\A_k$ is isomorphic to $\ker(p) \rtimes \SAut(\F_k)$, where $\SAut(\F_k)$ acts on $\ker(p)$ by conjugation. 
Let $\psi \in \ker(p)$, then $\psi$ fixes all $x_1,\ldots,x_k$ and since it is generated by $C$-generators and $N$-generators, to $x_0$ we can multiply only elements from $\F_k$ on the right and left. Thus there exist unique elements $l_{\psi}$ and $w_{\psi}$ in $\F_k$ such that:  
$$\psi(x_0,x_1,\ldots,x_k) = (l_{\psi}x_0r_{\psi},x_1,\ldots,x_k).$$

The map $j(\psi) = (l_\psi^{-1},r_{\psi})$ is an isomorphism between $\ker(p)$ and $\F_k^2$. Indeed, it is injective and since we have all generators $L^{\pm}_{0i}$ and $R^{\pm}_{0i}$, we can obtain every element in $\F_k^2$. Moreover, since
$$
\psi_1\psi_2(x_0) = \psi_1(l_{\psi_2}x_0r_{\psi_2}) = l_{\psi_2}l_{\psi_1}x_0r_{\psi_1}r_{\psi_2},
$$

we have that $j(\psi_1\psi_2) = (l_{\psi_1}^{-1}l_{\psi_2}^{-1},r_{\psi_1}r_{\psi_2}) = j(\psi_1)j(\psi_2)$, thus $j$ is a homomorphism. It follows that $\A_k$ is isomorphic to $\F_k^2 \rtimes \SAut(\F_k)$.

Now we compute how the conjugation action of $\SAut(\F_k)$ on $\ker(p)$ looks like after identifying $\ker(p)$ with $\F_k^2$ via the isomorphism $j$. Let $\phi \in \SAut(\F_k) < \A_k$ and $\psi \in \ker(p)$ be such that $j(\psi) = (l^{-1},r)$. Note that $\phi(x_0) = x_0$. We have 

\begin{align*}
\phi \psi \phi^{-1} (x_0) & =\phi \psi (x_0)\\
& = \phi (lx_0r)\\
&=\phi(l)x_0\phi(r).
\end{align*}

Thus $j(\phi \psi \phi^{-1}) = (\phi(l^{-1}),\phi(r))$ and the action of $\SAut(\F_k)$ on $\F_k^2$ is the direct product of the canonical action.

\end{proof}

\end{document}
